Benchmark scores for network intrusion detectors hide problems that appear only in deployment. We audit a decision tree on 54 NFStream flow features from labels to live blocking on an emulated Open vSwitch gateway, and condense the audit into four checks. \emph{Labels:} 91--98.5\% of the flows in the CIC-IoT2023 web, brute-force and MITM captures are background, i.e.\ benign traffic of other devices. Trained on file-name labels, the tree flags 71--88\% of the background flows in an attack capture it never saw, against 0.5--5.4\% with evidence-based labels; on IoT-23, with the authors' Zeek labels, it flags 23--47\% of unseen benign devices' flows, against 0--17.5\%. \emph{Splits:} on CIC-IoT2023 a random split hides that brute-force recall falls from 1.00 to 0.05 when the attack moves from SSH passwords to RTSP stream paths; this case (three host pairs) is most of the random--temporal gap. On IoT-23 the gap comes from classes missing from one half of the split. \emph{Live:} matching live to offline flows one by one exposed a forwarding pitfall, idle flow meters and an NFStream export-ordering defect that raised a median verdict delay to 82.5\,s, and showed slow replay splitting flows under unchanged macro-F1. \emph{Response:} blocking a host pair on its first attack verdict caused 60 false blocks in ten minutes of traffic; requiring 10 verdicts caused none but missed the web attackers and the Mirai bots, whose floods form too few flows. Code, attacker evidence and result files are public.
